\chapter{Introdução}
\label{cap:introducao}

\section{O problema de Pesquisa}\label{sec:citacoes}

    Historicamente, o ensino tem sido marcado pela divisão do conhecimento em disciplinas isoladas. Isso cria uma grande distância entre o que é estudado na universidade e os desafios que o mercado de trabalho realmente exige, dificultando que os alunos percebam a utilidade prática daquilo que estão aprendendo. Oliveira (2012) explica que, por muito tempo, a educação se deparou com um grande leque de conhecimentos e de teorias valiosas, porém trabalhadas de forma desconectada. Esse modelo tradicional, muitas vezes chamado de "ensino em caixinhas", faz com que o estudante sinta que a teoria acadêmica não tem relação direta com o seu dia a dia profissional.

    Para que o desenvolvimento científico e tecnológico avance de verdade e gere um impacto positivo, é necessário superar essa separação dos saberes. Ribeiro et al. (2025) destacam que problemas reais da sociedade contemporânea raramente respeitam os limites de uma única disciplina; eles são complexos e exigem respostas que combinem várias áreas ao mesmo tempo. Nessa mesma linha, Bentes (2024) reforça que as instituições de ensino precisam modernizar seus currículos, conectando a pesquisa e a educação com a inovação, para que a universidade consiga atuar de forma prática e direta nas necessidades da comunidade. Sem essa mudança de postura, o aprendizado corre o risco de permanecer apenas na teoria.

    Diante desse cenário, onde a separação das matérias ainda é um obstáculo para a formação de bons profissionais, define-se o seguinte problema de pesquisa: \textbf{De que forma a integração multidisciplinar e interdisciplinar, aliada a metodologias ativas de ensino que estimulam o protagonismo do aluno, pode superar a fragmentação do conhecimento e impulsionar o desenvolvimento de soluções mais assertivas e inovadoras para problemas complexos?} A busca por essa resposta é o que guia o presente trabalho, exigindo uma nova visão sobre como o processo de ensino em Engenharia de Software deve ser conduzido para garantir resultados funcionais.
    
\section{A hipótese}

    A hipótese central deste trabalho é que o protagonismo do aluno proporciona a qualidade no aprendizado quando ele é inserido em um ambiente que utiliza metodologias ativas de ensino. Quando o aluno deixa de ser um mero espectador e passa a ser o construtor do próprio projeto, como no caso da Aprendizagem Baseada em Problemas (PBL), o nível de entendimento e retenção do conteúdo aumenta consideravelmente. Ao assumir a responsabilidade pela criação da solução, o estudante se sente motivado a buscar respostas em diferentes disciplinas ao mesmo tempo, quebrando de vez a barreira do ensino isolado.

    Ao substituir o isolamento das matérias por uma abordagem multidisciplinar e interdisciplinar, o processo educacional passa a refletir as reais dificuldades e a complexidade do mundo corporativo. Acredita-se que, ao permitir que o aluno interaja com múltiplas áreas do conhecimento de forma integrada (formando uma verdadeira "rede de saberes"), ele desenvolve uma maior capacidade de análise, pensamento crítico, criatividade e habilidades para trabalhar em equipe. Oliveira (2012) defende que o pensamento integrado ajuda o estudante a relacionar os conteúdos de forma sistêmica, tornando o aprendizado algo natural e útil.

    Consequentemente, essa união entre disciplinas atua como um motor essencial para o desenvolvimento de software e tecnologias. Ao trabalhar de forma colaborativa, utilizando métodos práticos de gestão e comunicação da indústria, o estudante é capacitado a desenhar e executar projetos reais. Cruz et al. (2024) evidenciam que equipes engajadas em propósitos comuns e que utilizam métodos integrados têm um desempenho superior na criação de produtos. Dessa forma, a vivência acadêmica culmina na concepção de softwares e soluções muito mais assertivas, práticas e com grande capacidade de inovação para o mercado.

\section{Objetivos}

    \textbf{Objetivo Geral:} Analisar como a integração multidisciplinar e interdisciplinar no ambiente educacional, impulsionada por metodologias ativas de ensino e pelo protagonismo do aluno, contribui para o desenvolvimento científico e proporciona a criação de soluções mais assertivas para os desafios complexos da sociedade e do mercado.

    \textbf{Objetivos Específicos:} Para que o objetivo geral desta pesquisa seja atingido de maneira estruturada, foram definidos os seguintes passos práticos, focados na avaliação das disciplinas de Projeto Integrador (PIES) do curso de Engenharia de Software na UFC (Universidade Federal do Ceará):

    \begin{itemize}
        \item \textbf{Identificar os principais desafios e barreiras} que os estudantes enfrentam em relação ao ensino dividido em disciplinas isoladas e à falta de conexão prática no início do seu aprendizado técnico.
        \item \textbf{Mapear como os conteúdos das disciplinas concomitantes} de cada semestre se unem e oferecem o suporte técnico e de negócios necessário para o bom andamento das disciplinas de Projeto Integrador (PIES 1, PIES 2 e PIES 3).
        \item \textbf{Desenvolver e aplicar formulários de pesquisa} para os alunos das disciplinas de PIES, a fim de coletar dados concretos e métricas sobre o que eles acham a respeito de usar várias áreas do conhecimento em um mesmo projeto.
        \item \textbf{Avaliar o impacto das metodologias ativas e do protagonismo estudantil} na qualidade do aprendizado, verificando se a atitude ativa e participativa do aluno ajuda a evitar o retrabalho e a acelerar o tempo de desenvolvimento do software.
        \item \textbf{Analisar, por meio dos resultados coletados,} se a união de diferentes saberes e matérias ao longo do curso resulta na entrega final de softwares e soluções tecnológicas mais precisas, funcionais e inovadoras para o cliente final.
    \end{itemize}